\documentclass[11pt]{article}
\usepackage[margin=1in]{geometry}
\usepackage{amsmath,amssymb,mathtools}
\usepackage[hidelinks]{hyperref}

\usepackage{amsmath,amssymb}
\usepackage{algorithm}
\usepackage{algpseudocode}

\newcommand{\Pa}{\operatorname{Pa}}
\newcommand{\Ch}{\operatorname{Ch}}
\newcommand{\MB}{\operatorname{MB}}
\newcommand{\indep}{\perp\!\!\!\perp}

\title{Structural Causal Models and Markov Blankets}
\author{}
\date{}

\begin{document}
\maketitle

\section{Structural causal models}

Let $V=\{V_1,\ldots,V_n\}$ be the endogenous variables in an acyclic
structural causal model (SCM), and let $\Pa(V_i)$ denote the parents of
$V_i$ in its directed acyclic graph (DAG). The SCM consists of structural
equations
\[
  V_i=f_i\bigl(\Pa(V_i),U_i\bigr), \qquad i=1,\ldots,n,
\]
where $U_i$ are exogenous variables. If the exogenous variables are mutually
independent, the observational distribution obeys the causal Markov
factorization
\[
  P(V_1,\ldots,V_n)=\prod_{i=1}^n P\bigl(V_i\mid\Pa(V_i)\bigr).
\]

Edges are ordinarily structural features of the DAG rather than random
variables. If graph structure is uncertain and represented by a random graph
$G$, a joint model is
\[
  P(G,V)=P(G)\prod_{i=1}^n P\bigl(V_i\mid\Pa_G(V_i),G\bigr).
\]

\section{Formal definition of a Markov blanket}

For a variable $X\in V$, a set $B\subseteq V\setminus\{X\}$ is a
\emph{Markov blanket} of $X$ with respect to $P$ if
\[
  X\indep V\setminus\bigl(B\cup\{X\}\bigr)\mid B.
\]
Equivalently,
\[
  P\bigl(X\mid V\setminus\{X\}\bigr)=P(X\mid B).
\]
Thus, conditional on $B$, every variable outside $B$ carries no further
information about $X$.

A \emph{Markov boundary} is a minimal Markov blanket:
\[
  \MB(X)\in\arg\min_B
  \left\{B\subseteq V\setminus\{X\}:\;
  X\indep V\setminus\bigl(B\cup\{X\}\bigr)\mid B\right\},
\]
where minimality means that no member can be removed while preserving the
conditional-independence property. Under standard positivity/intersection
conditions, this boundary is unique.

For a distribution faithful to a DAG, the Markov boundary of $X$ is
\[
  \MB(X)=\Pa(X)\cup\Ch(X)\cup
  \left(\bigcup_{Y\in\Ch(X)}\Pa(Y)\setminus\{X\}\right).
\]
It therefore contains the parents of $X$, the children of $X$, and the other
parents of those children (often called spouses of $X$).

\section{Combining all Markov blankets}

The set of local Markov blankets is
\[
  \mathcal{B}=\left\{\MB(X_i): i=1,\ldots,n\right\}.
\]
It defines an undirected Markov-blanket graph
\[
  G_{\mathrm{MB}}=(V,E_{\mathrm{MB}}),
\]
where a symmetric edge rule is
\[
  \{X_i,X_j\}\in E_{\mathrm{MB}}
  \quad\Longleftrightarrow\quad
  X_j\in\MB(X_i)\ \text{or}\ X_i\in\MB(X_j).
\]
For an exactly known DAG, this is its moral graph: retain parent--child
adjacencies, connect co-parents of each child, and remove arrow directions.
Consequently, $G_{\mathrm{MB}}$ does not in general identify causal directions.

\section{Probability of a Markov-blanket graph}

A graph has a probability only when graph structure is modeled as uncertain.
Given observed data $D$, the posterior probability of a candidate
Markov-blanket graph is
\[
  P(G_{\mathrm{MB}}\mid D)=
  \frac{P(D\mid G_{\mathrm{MB}})P(G_{\mathrm{MB}})}
       {\sum_{G'}P(D\mid G')P(G')}.
\]
Here $P(G_{\mathrm{MB}})$ is a graph prior, often chosen to favor sparse
graphs, and $P(D\mid G_{\mathrm{MB}})$ is the marginal likelihood.

If $G_{\mathrm{MB}}$ is obtained by moralizing an uncertain causal DAG $G$,
then
\[
  P(G_{\mathrm{MB}}\mid D)=
  \sum_{G:\,\operatorname{moral}(G)=G_{\mathrm{MB}}}P(G\mid D).
\]
Several DAGs can contribute to the same Markov-blanket graph.

\section{Estimating a Markov blanket from data}

For a target $X$, estimate the smallest set $S$ satisfying
\[
  X\indep V\setminus\bigl(\{X\}\cup S\bigr)\mid S.
\]
The resulting estimate is $\widehat{\MB}(X)=S$. A constraint-based approach
iteratively adds variables $Y$ that remain conditionally dependent on $X$,
\[
  X\not\indep Y\mid S,
\]
and then removes a selected variable $Y\in S$ if
\[
  X\indep Y\mid S\setminus\{Y\}.
\]

For linear-Gaussian variables, conditional independence can be tested with
partial correlation:
\[
  \rho_{XY\mid S}\approx0\quad\Longrightarrow\quad X\indep Y\mid S.
\]
For discrete variables, use conditional mutual information:
\[
  I(X;Y\mid S)=\sum_{x,y,s}P(x,y,s)
  \log\!\left(\frac{P(x,y\mid s)}{P(x\mid s)P(y\mid s)}\right).
\]
Small values of $I(X;Y\mid S)$ support conditional independence. In
high-dimensional linear settings, neighborhood selection may use the Lasso:
\[
  \widehat{\beta}=\arg\min_{\beta}\left[
  \frac{1}{2n}\lVert X-\mathbf{X}_{-X}\beta\rVert_2^2+
  \lambda\lVert\beta\rVert_1\right],
  \qquad
  \widehat{\MB}(X)=\{Y:\widehat{\beta}_Y\ne0\}.
\]



\begin{algorithm}[H]
\caption{Bayesian Sequential Reconstruction of a Structural Causal Model}
\label{alg:bayesian_scm}
\begin{algorithmic}[1]

\State Let
\[
G=(V,E), \qquad
D=\text{observational data}, \qquad
D^{do}=\text{interventional data}.
\]

\Repeat

%---------------------------------------------------------
\State \textbf{Select a vertex}

Choose a vertex
\[
V_i\in V,
\]
for example,
\[
V_i
=
\arg\max_{V}
H(B(V)\mid D),
\]
where
\[
H(B)
=
-\sum_B
P(B\mid D)\log P(B\mid D).
\]

%---------------------------------------------------------
\State \textbf{Apply do-intervention and propagate a signal}

Perform the intervention
\[
do(V_i=x^*)
\]
and estimate
\[
P(X\mid do(V_i)).
\]

Forward diffusion is computed by
\[
u(t)=e^{-tL}u_0,
\qquad
u_0=e_i,
\]
where
\[
L=D-A
\]
is the graph Laplacian.

Backward propagation is obtained from the adjoint equation
\[
u^*(t)=e^{tL}u_T.
\]

%---------------------------------------------------------
\State \textbf{Sample candidate edges}

For every candidate vertex $V_j$ compute
\[
P(E_{ij}=1)
=
\sigma(s_{ij}),
\]
where
\[
s_{ij}
=
I(X_i;X_j)
+
\alpha
D_{KL}
\!\left(
P(X_j\mid do(X_i))
\|
P(X_j)
\right).
\]

Sample
\[
E_{ij}
\sim
P(E_{ij}\mid D,D^{do}).
\]

%---------------------------------------------------------
\State \textbf{Correct the Markov blanket}

Update
\[
B(V_i)
=
Pa(V_i)
\cup
Ch(V_i)
\cup
Sp(V_i).
\]

Estimate the posterior
\[
P(B(V_i)\mid D)
\propto
P(D\mid B(V_i))
P(B(V_i)).
\]

Choose
\[
B^*(V_i)
=
\arg\max_B
P(B\mid D).
\]

%---------------------------------------------------------
\State \textbf{Estimate local parameters and hyperparameters}

Estimate
\[
(\theta_i,\lambda_i)
=
\arg\max
P(D\mid\theta_i,B(V_i))
P(\theta_i\mid\lambda_i)
P(\lambda_i).
\]

%---------------------------------------------------------
\State \textbf{Insert the blanket into the graph}

Update
\[
G_{k+1}
=
G_k
\cup
B(V_i),
\]
or equivalently
\[
E_{k+1}
=
E_k
\cup
E(B(V_i)).
\]

%---------------------------------------------------------
\State \textbf{Fit global parameters and hyperparameters}

Estimate
\[
(\Theta,\Lambda)
=
\arg\max
P(D\mid G,\Theta)
P(\Theta\mid\Lambda)
P(\Lambda),
\]
where
\[
\Theta
=
(\theta_1,\ldots,\theta_n),
\qquad
\Lambda
=
(\lambda_1,\ldots,\lambda_n).
\]

For Gaussian SCMs this is equivalent to minimizing
\[
\Theta
=
\arg\min_{\Theta}
\left[
-\log
P(D\mid G,\Theta)
+
\sum_i
\lambda_i
\|\theta_i\|^2
\right].
\]

%---------------------------------------------------------
\State \textbf{Compute model evidence}

Compute the Bayesian evidence
\[
P(D\mid G)
=
\int
P(D\mid\Theta,G)
P(\Theta\mid G)
\,d\Theta.
\]

The Laplace approximation is
\[
\log P(D\mid G)
\approx
\log
P(D\mid\hat{\Theta},G)
+
\log
P(\hat{\Theta})
-
\frac12
\log|H|,
\]
where
\[
H
=
-\nabla^2
\log
P(\Theta\mid D).
\]

\Until
\[
\left|
\log P(D\mid G_{k+1})
-
\log P(D\mid G_k)
\right|
<
\varepsilon.
\]

\end{algorithmic}
\end{algorithm}

\section{Sources}
\begin{itemize}
  \item J. Pearl, ``The Foundations of Causal Inference,'' \emph{Sociological
  Methodology}, 2010. \url{https://doi.org/10.1111/j.1467-9531.2010.01228.x}
  \item J. Peters, D. Janzing, and B. Sch\"olkopf, \emph{Elements of Causal
  Inference}, MIT Press, 2017. \url{https://web.math.ku.dk/~peters/jonas_files/ElementsOfCausalInference.pdf}
\end{itemize}

\end{document}
